\newcommand{\CalibrationStatus}{Exact Lean theorem \code{Q222 F} iff \code{BProv ZFCax_s [] Form.fBot}; all 110 modules freshly compiled.}
\newcommand{\CalibrationDiscussion}{The four supplied modules finish this argument as an exact Lean theorem:
\begin{quote}\small\code{Quotient222.Q222 Code254.F} \(\longleftrightarrow\)\\\code{BProv ZFCax_s [] Form.fBot}.\end{quote}
The upper adapter uses \code{KUpper222.bundles_in_L}; the empty set and its successor supply distinct elements of the constructible interpretation. The lower adapter reuses the short mathematical argument from \code{KEndpoint.c222_of_zf_refutation}, followed by \code{Q222.base}. It imports \code{KLowerQHyp} directly, avoiding unrelated tape-machine endpoints. All 110 modules were compiled afresh under Lean 4.32.0, with no cached project proof objects: 106 upstream modules match the pinned Git blobs and four new modules total 244 lines including comments and blanks. The final axiom report is exactly \code{propext}, \code{Classical.choice}, and \code{Quot.sound}; there is no added consistency axiom. This is refutational equivalence, not equality of all theorem sets, and it does not itself establish a new binary-machine endpoint.}
\newcommand{\RocqDiscussion}{A direct Coq companion compiled under Coq 8.18.0 and passed an independent recursive \code{coqchk}. Its \code{rho8_all} assumption audit is closed under the global context. Its 24,352-byte, 440-line source is larger in bytes than the current 22,673-byte, 1,169-line reflected Rocq certificate. It offers the same direct proof structure without the reflection interpreter and 735-instruction table; it is not a claimed Rocq byte reduction. The two imported numerical modules match the pinned upstream Git blobs without edits. This independent check used Coq 8.18.0 rather than the full endpoint's Rocq 9.1.0 toolchain.}
